\pagestyle{empty}
\tight
\begin{tblr}{width=\textwidth,colspec={|Q[c,m,1.9cm]|X[l,m]|},hlines,vlines,hspan=minimal,row{1}={bg=headblue},rowsep=1pt,colsep=3pt}
\SetCell{c,m}\hfil\logo\hfil & \textbf{POLITEKNIK NEGERI MALANG}\par\textbf{JURUSAN TEKNOLOGI INFORMASI}\par\textbf{PROGRAM STUDI: D4 TEKNIK INFORMATIKA} \\
\SetCell[c=2]{c} \textbf{RENCANA PEMBELAJARAN SEMESTER (RPS)} & \\
\end{tblr}
\begin{longtblr}{width=\textwidth,colspec={|Q[l,m,1.9cm]|X[0.85,l,m]|X[1.45,l,m]|X[0.75,l,m]|X[1.15,l,m]|},hlines,vlines,hspan=minimal,rowsep=1pt,colsep=2pt,row{1,3}={bg=headgray,font=\bfseries}}
MATA KULIAH & KODE & BOBOT (SKS)/jam & SEMESTER & TGL. PENYUSUNAN \\
Administrasi dan Keamanan Jaringan & RTI255008 & 3 SKS/6 jam & 5 & 15 Agustus 2025 \\
OTORISASI & Dosen Pengembang RPS & Koordinator RMK & \SetCell[c=2]{l} Ka PRODI &  \\
 & \begin{enumerate}\item Ade Ismail\item Yuri Ariyanto\item Sofyan Noor Arief\end{enumerate} &  & \SetCell[c=2]{l}{\vspace{8mm}\par Dr. Ely Setyo Astuti, S.T., M.T.} &  \\
\SetCell[r=14]{l} Capaian Pembelajaran (CP) & \SetCell[c=4]{l,bg=headgray,font=\bfseries} Capaian Pembelajaran Lulusan Program Studi (CPL-Prodi) &  &  &  \\
 & CPL05 & \SetCell[c=3]{l} Solusi teknologi komputasi dan infrastruktur digital sesuai kebutuhan. &  &  \\
 & \SetCell[c=4]{l,bg=headgray,font=\bfseries} Capaian Pembelajaran mata kuliah (CPL-MK) &  &  &  \\
 & CPMK0501 & \SetCell[c=3]{l} Mahasiswa mampu merancang dan menerapkan kebijakan serta prosedur keamanan infrastruktur jaringan dan cloud. &  &  \\
 & CPMK0502 & \SetCell[c=3]{l} Mahasiswa mampu menganalisis isu keamanan pada sistem operasi, jaringan, dan cloud untuk menentukan strategi mitigasi yang tepat. &  &  \\
 & CPMK0505 & \SetCell[c=3]{l} Mahasiswa mampu mengkonfigurasi, mengelola, dan mengoptimalkan mekanisme keamanan jaringan sesuai standar dan best practice. &  &  \\
 & CPMK0504 & \SetCell[c=3]{l} Mampu membangun dan mengelola infrastruktur jaringan komputer dengan memilih arsitektur, protokol, dan teknik konfigurasi yang sesuai kebutuhan sistem. &  &  \\
 & CPMK1002 & \SetCell[c=3]{l} Mampu merancang, menerapkan, dan mengevaluasi mekanisme keamanan jaringan untuk melindungi data, layanan, dan sumber daya jaringan dengan mempertimbangkan aspek teknis, operasional, dan kebijakan organisasi. &  &  \\
 & CPMK1003 & \SetCell[c=3]{l} Mampu menganalisis permasalahan kompleks pada infrastruktur dan layanan jaringan komputer untuk merancang, mengonfigurasi, dan mengelola solusi jaringan on-premise dan cloud yang andal, aman, dan scalable sesuai kebutuhan organisasi. &  &  \\
 & \SetCell[c=4]{l,bg=headgray,font=\bfseries} Kemampuan akhir tiap tahapan belajar (Sub-CPMK) &  &  &  \\
 & SCPMK0501-04401 & \SetCell[c=3]{l} Mahasiswa mampu menganalisis dan merancang kebijakan keamanan infrastruktur jaringan termasuk desain topologi aman dan pemilihan gateway internet. &  &  \\
 & SCPMK0502-04401 & \SetCell[c=3]{l} Mahasiswa mampu mengidentifikasi dan menganalisis ancaman, kerentanan, serta dampak keamanan pada sistem dan jaringan komputer. &  &  \\
 & SCPMK0505-04401 & \SetCell[c=3]{l} Mahasiswa mampu mengkonfigurasi dan menguji mekanisme keamanan jaringan meliputi NAT, firewall, administrasi akses, dan monitoring SIEM. &  &  \\
 & SCPMK0505-04402 & \SetCell[c=3]{l} Mahasiswa mampu melakukan optimasi kinerja jaringan, troubleshooting, benchmarking, dan mengintegrasikan SIEM untuk perencanaan kapasitas. &  &  \\
 & SCPMK0504-04403 & \SetCell[c=3]{l} Mampu membangun infrastruktur jaringan aman melalui desain arsitektur, segmentasi, dan konfigurasi gateway. &  &  \\
 & SCPMK1002-04404 & \SetCell[c=3]{l} Mampu merancang dan mengevaluasi mekanisme perlindungan data serta sumber daya menggunakan monitoring dan logging. &  &  \\
 & SCPMK1003-04405 & \SetCell[c=3]{l} Mampu menganalisis permasalahan kompleks dan mengelola solusi optimasi jaringan yang aman dan scalable. &  &  \\
\SetCell[r=2]{l} Penilaian & \SetCell[c=4]{l} *Nilai CPMK didapat dari agregasi nilai SUB-CPMK yang dibebankan &  &  &  \\
 & \SetCell[c=4]{l}{\tiny\begin{tblr}{width=\linewidth,colspec={|Q[c,m,0.1\linewidth]|Q[c,m,0.16\linewidth]|X[l,m]|X[l,m]|X[l,m]|X[l,m]|X[l,m]|Q[c,m,0.08\linewidth]|},hlines,vlines,hspan=minimal,rowsep=1pt,colsep=0.7pt,row{1,2}={bg=headgray,font=\bfseries}}
Kode CPMK & Kode SCPMK & \SetCell[c=5]{c} Bobot per Bentuk Penilaian & & & & & Total Bobot \\
 & & Tugas & Kuis & UTS & PBL & UAS & \\
\SetCell[r=1]{c} CPMK0501 & SCPMK0501-04401 & 6\%: kebijakan \& desain topologi & 4\%: konfigurasi gateway & 10\%: konsep \& infrastruktur & & & 20\% \\
\SetCell[r=1]{c} CPMK0502 & SCPMK0502-04401 & 4\%: ancaman \& kerentanan & & 10\%: identifikasi ancaman & 13\%: simulasi desain keamanan & & 27\% \\
\SetCell[r=2]{c} CPMK0505 & SCPMK0505-04401 & 6\%: administrasi \& monitoring & 6\%: mekanisme keamanan & & 13\%: simulasi administrasi jaringan & & 25\% \\
 & SCPMK0505-04402 & 8\%: optimasi \& troubleshooting & & & & 20\%: kompetensi akhir & 28\% \\
CPMK0504 & SCPMK0504-04403 & 0\% & 0\% & 0\% & 0\% & 0\% & terintegrasi dengan penilaian Minggu 3, 6 \\
CPMK1002 & SCPMK1002-04404 & 0\% & 0\% & 0\% & 0\% & 0\% & terintegrasi dengan penilaian Minggu 11 \\
CPMK1003 & SCPMK1003-04405 & 0\% & 0\% & 0\% & 0\% & 0\% & terintegrasi dengan penilaian Minggu 14--15 \\
\SetCell[c=7]{r}\textbf{Total Per Penilaian} & & & & & & & \textbf{100\%} \\
\end{tblr}} &  &  &  \\
Diskripsi Singkat Mata Kuliah & \SetCell[c=4]{l} Mata kuliah ini membekali mahasiswa dengan pengetahuan dan keterampilan komprehensif untuk menjadi administrator jaringan yang kompeten dalam mengelola keamanan jaringan komputer di organisasi. Materi meliputi desain sistem keamanan, instalasi dan konfigurasi gateway internet, administrasi sistem jaringan, optimasi kinerja jaringan, dan integrasi SIEM, sesuai Standar Kompetensi Kerja Nasional Indonesia (SKKNI) bidang Jaringan Komputer dan Sistem Administrasi. &  &  &  \\
Bahan Kajian / Materi Pembelajaran & \SetCell[c=4]{l}\begin{enumerate}\item Pengenalan infrastruktur \& keamanan jaringan, CIA Triad, dan penilaian risiko\item Identifikasi ancaman siber dan kerentanan sistem jaringan\item Desain topologi jaringan aman: segmentasi, VLAN, DMZ\item Simulasi desain keamanan jaringan\item Konsep dan pemilihan gateway internet\item Instalasi dan konfigurasi dasar gateway internet\item Konfigurasi NAT dan firewall pada gateway\item Administrasi sistem jaringan: AAA, RADIUS/TACACS+\item Manajemen pengguna, RBAC, dan kontrol akses\item Monitoring jaringan dan dasar logging/SIEM\item Simulasi administrasi dan troubleshooting jaringan\item Identifikasi bottleneck dan optimasi kinerja jaringan\item Benchmarking, QoS, dan traffic shaping\item Simulasi optimasi kinerja dan integrasi SIEM\end{enumerate} &  &  &  \\
\SetCell[r=2]{l} Pustaka & \SetCell[c=4]{l} \textbf{Utama:}\begin{enumerate}\item Standar Kompetensi Kerja Nasional Indonesia (SKKNI) Sektor Komunikasi dan Informasi, Bidang Jaringan Komputer dan Sistem Administrasi.\item Stallings, W. Cryptography and Network Security: Principles and Practice. Pearson.\item Cisco Networking Academy. CCNA Security Course Booklet. Cisco Press.\item NIST Special Publications (SP 800 Series) terkait Keamanan Jaringan.\end{enumerate} &  &  &  \\
 & \SetCell[c=4]{l} \textbf{Pendukung:} Materi ajar/modul dari dosen/tim pengajar; dokumentasi resmi GNS3 dan VirtualBox. &  &  &  \\
Media Pembelajaran & \SetCell[c=2]{l} \textbf{Software:} GNS3, VirtualBox. &  & \SetCell[c=2]{l} \textbf{Hardware:} Komputer/laptop, proyektor. &  \\
Nama Dosen Pengampu & \SetCell[c=4]{l}\begin{enumerate}\item Ade Ismail\item Yuri Ariyanto\item Sofyan Noor Arief\item Yusriel A\end{enumerate} &  &  &  \\
Matakuliah Syarat & \SetCell[c=4]{l} Sistem Operasi; Jaringan Komputer &  &  &  \\
\end{longtblr}
\begin{longtblr}{width=\textwidth,colspec={|Q[c,m,0.06\textwidth]|X[1.35,l,m]|X[1.08,l,m]|X[1.16,l,m]|X[0.7,l,m]|X[1.24,l,m]|X[1.06,l,m]|X[0.98,l,m]|Q[c,m,0.05\textwidth]|},hlines,vlines,hspan=minimal,rowhead=2,row{1,2}={bg=headgreen,font=\bfseries},rowsep=1pt,colsep=1.5pt}
Mgg.\par Ke & Kemampuan Akhir Yang Direncanakan (Sub-CP-MK) & Bahan kajian (Materi Pembelajaran) & Bentuk dan Metode Pembelajaran & Estimasi Waktu & Pengalaman Belajar Mahasiswa & Kriteria \& Bentuk Penilaian & Indikator Penilaian & Bobot Penilaian (\%) \\
(1) & (2) & (3) & (4) & (5) & (6) & (7) & (8) & (9) \\
1 & SCPMK0502-04401 & Pengantar infrastruktur \& keamanan jaringan. & Modalitas: Blended Learning. Bentuk: Luring/Daring. Metode: Ceramah, diskusi interaktif, studi kasus. & 1 x 3 x 50' tatap muka; 1 x 3 x 50' tugas/praktik mandiri. & Mahasiswa mengerjakan aktivitas terarah untuk pengantar infrastruktur dan keamanan jaringan serta menunjukkan evidence hasil belajar. & Bentuk penilaian: Tugas. Mengacu rubrik RTM. & Ketepatan konsep; kualitas artefak/praktik; validasi dan dokumentasi. & 2 \\
2 & SCPMK0502-04401 & Identifikasi ancaman \& kerentanan jaringan. & Modalitas: Blended Learning. Bentuk: Luring/Daring. Metode: Ceramah, diskusi, analisis kasus ancaman. & 1 x 3 x 50' tatap muka; 1 x 3 x 50' tugas/praktik mandiri. & Mahasiswa mengerjakan aktivitas terarah untuk identifikasi ancaman dan kerentanan serta menunjukkan evidence hasil belajar. & Bentuk penilaian: Tugas. Mengacu rubrik RTM. & Ketepatan konsep; kualitas artefak/praktik; validasi dan dokumentasi. & 2 \\
3 & SCPMK0501-04401, SCPMK0504-04403 & Desain topologi jaringan aman. & Modalitas: Blended Learning. Bentuk: Luring/Daring. Metode: Ceramah, diskusi kelompok, latihan desain konseptual. & 1 x 3 x 50' tatap muka; 1 x 3 x 50' tugas/praktik mandiri. & Mahasiswa mengerjakan aktivitas terarah untuk desain topologi jaringan aman serta menunjukkan evidence hasil belajar. & Bentuk penilaian: Tugas. Mengacu rubrik RTM. & Ketepatan konsep; kualitas artefak/praktik; validasi dan dokumentasi. & 2 \\
4 & SCPMK0501-04401, SCPMK0502-04401 & Simulasi desain keamanan jaringan. & Modalitas: Blended Learning. Bentuk: Luring/Daring. Metode: Praktikum/simulasi (lab), tugas proyek mini. & 1 x 3 x 50' tatap muka; 1 x 3 x 50' tugas/praktik mandiri. & Mahasiswa terlibat dalam proyek mini, melakukan simulasi desain keamanan, dan menyusun laporan dokumentasi. & Bentuk penilaian: PBL simulasi desain keamanan. Mengacu rubrik RTM. & Ketepatan konsep; kualitas artefak/praktik; validasi dan dokumentasi. & 13 \\
5 & SCPMK0501-04401 & Konsep \& pemilihan gateway internet. & Modalitas: Blended Learning. Bentuk: Luring/Daring. Metode: Ceramah, diskusi, analisis kebutuhan gateway. & 1 x 3 x 50' tatap muka; 1 x 3 x 50' tugas/praktik mandiri. & Mahasiswa mengerjakan aktivitas terarah untuk konsep dan pemilihan gateway internet serta menunjukkan evidence hasil belajar. & Bentuk penilaian: Tugas. Mengacu rubrik RTM. & Ketepatan konsep; kualitas artefak/praktik; validasi dan dokumentasi. & 2 \\
6 & SCPMK0505-04401, SCPMK0504-04403 & Instalasi dasar gateway internet. & Modalitas: Blended Learning. Bentuk: Luring/Daring. Metode: Ceramah, studi kasus, instalasi dan konfigurasi. & 1 x 3 x 50' tatap muka; 1 x 3 x 50' tugas/praktik mandiri. & Mahasiswa mengerjakan aktivitas terarah untuk instalasi dasar gateway internet serta menunjukkan evidence hasil belajar. & Bentuk penilaian: Tugas. Mengacu rubrik RTM. & Ketepatan konsep; kualitas artefak/praktik; validasi dan dokumentasi. & 2 \\
7 & SCPMK0501-04401, SCPMK0505-04401 & Kuis: konfigurasi NAT \& firewall pada gateway. & Modalitas: Blended Learning. Bentuk: Luring/Daring. Metode: Kuis tertulis. & 1 x 3 x 50' tatap muka; 1 x 3 x 50' tugas/praktik mandiri. & Mahasiswa mengerjakan kuis tertulis tentang konfigurasi NAT dan keamanan gateway. & Bentuk penilaian: Kuis konfigurasi dan keamanan jaringan. Mengacu rubrik RTM. & Ketepatan konsep; kualitas artefak/praktik; validasi dan dokumentasi. & 10 \\
8 & SCPMK0501-04401, SCPMK0502-04401 & UTS: infrastruktur \& keamanan jaringan. & Modalitas: Blended Learning. Bentuk: Luring/Daring. Metode: Ujian tertulis. & 1 x 3 x 50' tatap muka; 1 x 3 x 50' tugas/praktik mandiri. & Mahasiswa mengerjakan soal ujian tertulis komprehensif materi pertemuan 1--7. & Bentuk penilaian: UTS infrastruktur \& keamanan jaringan. Mengacu rubrik RTM. & Ketepatan konsep; kualitas artefak/praktik; validasi dan dokumentasi. & 20 \\
9 & SCPMK0505-04402 & Administrasi sistem jaringan. & Modalitas: Blended Learning. Bentuk: Luring/Daring. Metode: Ceramah, diskusi, latihan konfigurasi. & 1 x 3 x 50' tatap muka; 1 x 3 x 50' tugas/praktik mandiri. & Mahasiswa mengerjakan aktivitas terarah untuk administrasi sistem jaringan serta menunjukkan evidence hasil belajar. & Bentuk penilaian: Tugas. Mengacu rubrik RTM. & Ketepatan konsep; kualitas artefak/praktik; validasi dan dokumentasi. & 2 \\
10 & SCPMK0505-04402 & Manajemen pengguna \& kontrol akses. & Modalitas: Blended Learning. Bentuk: Luring/Daring. Metode: Ceramah, diskusi, latihan manajemen user. & 1 x 3 x 50' tatap muka; 1 x 3 x 50' tugas/praktik mandiri. & Mahasiswa mengerjakan aktivitas terarah untuk manajemen pengguna dan kontrol akses serta menunjukkan evidence hasil belajar. & Bentuk penilaian: Tugas. Mengacu rubrik RTM. & Ketepatan konsep; kualitas artefak/praktik; validasi dan dokumentasi. & 2 \\
11 & SCPMK0505-04402, SCPMK1002-04404 & Monitoring jaringan \& logging (SIEM). & Modalitas: Blended Learning. Bentuk: Luring/Daring. Metode: Ceramah, demo alat monitoring, diskusi. & 1 x 3 x 50' tatap muka; 1 x 3 x 50' tugas/praktik mandiri. & Mahasiswa mengerjakan aktivitas terarah untuk monitoring jaringan dan dasar SIEM serta menunjukkan evidence hasil belajar. & Bentuk penilaian: Tugas. Mengacu rubrik RTM. & Ketepatan konsep; kualitas artefak/praktik; validasi dan dokumentasi. & 2 \\
12 & SCPMK0505-04401, SCPMK0505-04402 & Simulasi administrasi \& troubleshooting jaringan. & Modalitas: Blended Learning. Bentuk: Luring/Daring. Metode: Praktikum/simulasi (lab), tugas proyek mini. & 1 x 3 x 50' tatap muka; 1 x 3 x 50' tugas/praktik mandiri. & Mahasiswa terlibat dalam proyek mini, melakukan simulasi administrasi dan troubleshooting, serta menganalisis skenario. & Bentuk penilaian: PBL simulasi administrasi \& troubleshooting. Mengacu rubrik RTM. & Ketepatan konsep; kualitas artefak/praktik; validasi dan dokumentasi. & 13 \\
13 & SCPMK0505-04402 & Identifikasi bottleneck kinerja jaringan. & Modalitas: Blended Learning. Bentuk: Luring/Daring. Metode: Ceramah, diskusi, analisis data kinerja. & 1 x 3 x 50' tatap muka; 1 x 3 x 50' tugas/praktik mandiri. & Mahasiswa mengerjakan aktivitas terarah untuk identifikasi bottleneck kinerja jaringan serta menunjukkan evidence hasil belajar. & Bentuk penilaian: Tugas. Mengacu rubrik RTM. & Ketepatan konsep; kualitas artefak/praktik; validasi dan dokumentasi. & 2 \\
14 & SCPMK0505-04402, SCPMK1003-04405 & Benchmarking \& tuning jaringan. & Modalitas: Blended Learning. Bentuk: Luring/Daring. Metode: Ceramah, studi kasus optimasi, diskusi. & 1 x 3 x 50' tatap muka; 1 x 3 x 50' tugas/praktik mandiri. & Mahasiswa mengerjakan aktivitas terarah untuk benchmarking dan tuning jaringan serta menunjukkan evidence hasil belajar. & Bentuk penilaian: Tugas. Mengacu rubrik RTM. & Ketepatan konsep; kualitas artefak/praktik; validasi dan dokumentasi. & 2 \\
15 & SCPMK0505-04402, SCPMK1003-04405 & Simulasi optimasi kinerja \& SIEM. & Modalitas: Blended Learning. Bentuk: Luring/Daring. Metode: Praktikum/simulasi, tugas mandiri. & 1 x 3 x 50' tatap muka; 1 x 3 x 50' tugas/praktik mandiri. & Mahasiswa mengerjakan aktivitas terarah untuk simulasi optimasi kinerja dan penggunaan SIEM serta menunjukkan evidence hasil belajar. & Bentuk penilaian: Tugas. Mengacu rubrik RTM. & Ketepatan konsep; kualitas artefak/praktik; validasi dan dokumentasi. & 2 \\
16 & SCPMK0505-04401, SCPMK0505-04402 & Refleksi \& portofolio kompetensi. & Modalitas: Blended Learning. Bentuk: Luring/Daring. Metode: Diskusi interaktif, presentasi ringkas, konsultasi portofolio. & 1 x 3 x 50' tatap muka; 1 x 3 x 50' tugas/praktik mandiri. & Mahasiswa merefleksikan pembelajaran, menyusun portofolio kompetensi, dan mendiskusikan keterkaitan teori-praktik. & Bentuk penilaian: Tugas. Mengacu rubrik RTM. & Ketepatan konsep; kualitas artefak/praktik; validasi dan dokumentasi. & 2 \\
17 & SCPMK0501-04401--SCPMK0505-04402 & UAS: kompetensi administrasi \& keamanan jaringan. & Modalitas: Blended Learning. Bentuk: Luring/Daring. Metode: Ujian tertulis atau presentasi proyek akhir. & 1 x 3 x 50' tatap muka; 1 x 3 x 50' tugas/praktik mandiri. & Mahasiswa mengerjakan soal ujian tulis atau mempresentasikan proyek akhir yang menunjukkan kompetensi SKKNI. & Bentuk penilaian: UAS kompetensi administrasi \& keamanan jaringan. Mengacu rubrik RTM. & Ketepatan konsep; kualitas artefak/praktik; validasi dan dokumentasi. & 20 \\
\end{longtblr}
